\documentclass[11pt,a4paper]{article}

\usepackage{iftex}
\ifPDFTeX
  \usepackage[utf8]{inputenc}
  \usepackage[T1]{fontenc}
  \usepackage{lmodern}
\else
  \usepackage{fontspec}
  \setmainfont{Liberation Serif}
  \setsansfont{Liberation Sans}
  \setmonofont{Liberation Mono}
\fi
\usepackage[czech]{babel}
\usepackage{amsmath,amssymb,amsthm,mathtools}
\usepackage[margin=25mm]{geometry}
\usepackage{microtype}
\usepackage[hidelinks]{hyperref}

\newtheorem{theorem}{Věta}
\newtheorem{lemma}[theorem]{Lemma}
\newtheorem{proposition}[theorem]{Tvrzení}
\theoremstyle{definition}
\newtheorem{definition}[theorem]{Definice}
\theoremstyle{remark}
\newtheorem{remark}[theorem]{Poznámka}

\newcommand{\Adm}{\operatorname{Adm}}
\newcommand{\Aobs}{A_{\mathrm{obs}}}

\title{\textbf{Principium relationis}\\[4pt]
\large totum relatum et imago accessibilis}
\author{Jan Mikulík}
\date{17. listopadu 2026}

\begin{document}
\maketitle

\begin{abstract}
Jméno je odečet. Dvě různé věci mohou dostat stejné jméno a přesto zůstat různé.
Vlastnost toto společné jméno dostane právě tehdy, když na jeho vlákně nerozlišuje.
Z odečtu proto nevyvozujeme vlastnost, která na vlákně není konstantní.
Akce bezpečná bez znalosti přesného předobrazu je průnik přes vlákno; prázdný průnik
se nenahrazuje domnělým souhlasem.
Text zapisuje věty ověřené v Leanu jako
\texttt{GoldenSphere.factors\_through\_iff},
\texttt{same\_name\_different\_things},
\texttt{splitting\_property\_does\_not\_factor},
\texttt{constant\_receives\_the\_name},
\texttt{A\_obs} a \texttt{obs\_mono}
v repozitáři \texttt{finite-Shannon-theorem}, commit \texttt{2e6b8e5}.
Bellovy stavy, ostrá nerovnost pro \(M_t\) a jakákoli identifikace s vědomím tu nejsou věty.
\end{abstract}

\section{Jinými slovy}

Henri Poincaré napsal, že matematika je umění dávat různým věcem stejné jméno.
V tomto textu se ta věta čte jen jako výrok o odečtu.

Nechť \(X\) je množina stavů, \(Y\) množina jmen a
\[
  O\colon X\to Y
\]
dostupný odečet. Jméno stavu \(x\) je \(O(x)\).
Vlákno jména \(y\in Y\) je
\[
  O^{-1}(y)=\{x\in X:O(x)=y\}.
\]
Společné jméno neznamená totožnost stavů. Znamená jen, že odečet je nerozlišil.

\begin{proposition}[Stejné jméno, různé věci]
Existuje odečet \(O\colon\{0,1\}\to\{\ast\}\) a stavy \(x,x'\) takové, že
\[
  O(x)=O(x')\qquad\text{a}\qquad x\neq x'.
\]
\end{proposition}

\begin{proof}
Položme \(O(0)=O(1)=\ast\), \(x=0\) a \(x'=1\).
\end{proof}

Latinské heslo \emph{totum relatum et imago accessibilis} tu znamená právě toto:
celek nese vztahy, které jednotlivý odečet nemusí uchovat.
Obraz je to, co je dostupné. Není to celý stav.

\section{Kdy vlastnost dostane jméno}

Nechť \(Z\) je množina hodnot a \(f\colon X\to Z\) vlastnost.
Řekneme, že \(f\) se faktoruje přes \(O\), existuje-li
\[
  g\colon O(X)\to Z
\]
takové, že \(f=g\circ O\).
Ekvivalentně: \(g\) je definováno na obrazu a
\[
  g\bigl(O(x)\bigr)=f(x)\qquad\text{pro každé }x\in X.
\]

\begin{theorem}[Rozhodnutelnost z obrazu]
\label{thm:factors}
Existuje \(g\colon O(X)\to Z\) s \(f=g\circ O\) právě tehdy, když
\begin{equation}
\label{eq:fibre}
  O(x)=O(x')\implies f(x)=f(x')
  \qquad\text{pro všechna }x,x'\in X.
\end{equation}
\end{theorem}

\begin{proof}
Nechť takové \(g\) existuje a \(O(x)=O(x')\).
Pak
\[
  f(x)=g\bigl(O(x)\bigr)=g\bigl(O(x')\bigr)=f(x').
\]
Obráceně nechť platí \eqref{eq:fibre}.
Pro \(y\in O(X)\) zvolme libovolný předobraz \(x_y\in O^{-1}(y)\) a položme
\(g(y)=f(x_y)\).
Je-li \(x'_y\) jiný předobraz, je \(O(x_y)=O(x'_y)\), tedy \(f(x_y)=f(x'_y)\).
Hodnota nezávisí na volbě. Pro každé \(x\) je \(x\) předobrazem \(O(x)\), takže
\(g(O(x))=f(x)\).
\end{proof}

Pracovní pravidlo je jedna věta.
Z odečtu nevyvozujeme vlastnost, která na jeho vlákně není konstantní.
Potřebujeme-li ji, rozšíříme odečet.
Už zaznamenanou událost tím zpětně neměníme.

\begin{proposition}[Vlastnost, která jméno rozdělí]
\label{prop:split}
Neexistuje \(g\colon\{\ast\}\to\{0,1\}\) takové, že
\[
  g\bigl(O(x)\bigr)=x\qquad\text{pro obě }x\in\{0,1\},
\]
kde \(O\) je konstantní odečet z předchozího tvrzení.
\end{proposition}

\begin{proposition}[Konstantní vlastnost jméno dostane]
\label{prop:const}
Nechť \(c\in Z\) a \(O\colon X\to\{\ast\}\) je konstantní.
Pak existuje \(g\colon\{\ast\}\to Z\) s \(g(O(x))=c\) pro každé \(x\in X\).
\end{proposition}

\begin{proof}[Důkaz obou tvrzení]
Identita na \(\{0,1\}\) nesplňuje \eqref{eq:fibre}, protože \(O(0)=O(1)\) a \(0\neq 1\).
Věta~\ref{thm:factors} tedy faktor zakazuje.
Konstantní \(f(x)=c\) podmínku \eqref{eq:fibre} splňuje, takže faktor existuje:
\(g(\ast)=c\).
\end{proof}

Společné jméno tedy není ztráta rozlišení uvnitř matematiky.
Je to přesná hranice toho, co smí nést totéž označení.
Věci zůstávají různé. Vlastnost, která je rozliší, jejich jméno nedostane.

\section{Akce bez znalosti předobrazu}

Data sféry jsou čtveřice
\[
  I\colon W\times X\to \mathcal{P}(\Omega),\qquad
  F\colon X\times\Omega\to X,\qquad
  C\colon N\times X\to \mathcal{P}(\Omega),\qquad
  K\subseteq X.
\]
\(I(w,x)\) jsou akce otevřené přáním, \(F(x,a)\) je následný stav,
\(C(i,x)\) je hranice souhlasu činitele \(i\) a \(K\) je množina péče.

\begin{definition}[Přípustná akce]
\[
  \Adm(w,x)
  =
  \bigl\{\,a\in\Omega :
    a\in I(w,x)
    \ \text{a}\ 
    \forall i\in N,\; a\in C(i,x)
  \,\bigr\}.
\]
Nabízený následný stav je
\[
  G(w,x)
  =
  \bigl\{\,F(x,a):
    a\in\Adm(w,x)
    \ \text{a}\ 
    F(x,a)\in K
  \,\bigr\}.
\]
\end{definition}

Každý prvek \(G(w,x)\) leží v \(K\).
Definice nemá klauzuli, která by chybějící souhlas nahradila obětí jiného činitele.
Existence prvku v \(G(w,x)\) je samostatný předpoklad: sféra svědka nevyrábí.

Při odečtu \(y\) jsou nerozlišené stavy právě vlákno \(O^{-1}(y)\).
Akci, která je bezpečná, ať je předobraz kterýkoli, bereme z průniku.

\begin{definition}[Bezpečné akce při odečtu]
\label{def:aobs}
\[
  \Aobs(w,y)
  =
  \bigl\{\,a\in\Omega :
    \forall x\in O^{-1}(y),\;
    a\in\Adm(w,x)
    \ \text{a}\ 
    F(x,a)\in K
  \,\bigr\}.
\]
\end{definition}

\begin{proposition}[Menší vlákno průnik nezmenší]
\label{prop:mono}
Nechť \(M'\subseteq M\subseteq X\). Pak
\begin{align*}
  &\bigl\{\,a :
    \forall x\in M,\;
    a\in\Adm(w,x)\ \text{a}\ F(x,a)\in K
  \,\bigr\}\\
  &\qquad\subseteq
  \bigl\{\,a :
    \forall x\in M',\;
    a\in\Adm(w,x)\ \text{a}\ F(x,a)\in K
  \,\bigr\}.
\end{align*}
\end{proposition}

\begin{proof}
Splňuje-li \(a\) podmínku pro každé \(x\in M\), splňuje ji i pro každé \(x\in M'\),
protože \(M'\subseteq M\).
\end{proof}

Nový odečet může vlákno zúžit a tím průnik zvětšit.
Může tedy uvolnit volbu, která předtím v průniku nebyla.
Nevymýšlí souhlas, který ve zbylých stavech chybí.
Je-li průnik prázdný, akce při tomto odečtu není.

\section{Co tu není větou}

Následující věci tento text nedokazuje a z předchozích vět neplynou.

\begin{enumerate}
\item Že přání existuje, že hranice souhlasu jsou ty správné,
      že \(K\) je neprázdná, nebo že \(G(w,x)\) je neprázdná.
\item Že Bellovy stavy \(\lvert\Phi_\pm\rangle\) jsou modelem vědomí.
      Shoda lokálních odečtů při různém společném stavu je příklad vlákna,
      ne důkaz fyzikální ontologie.
\item Že symetrie \(M_t(x)=M_t(-x)\) vynucuje nuly funkce \(\xi\) na kritické přímce.
      Postačující ostrá nerovnost \(M_t(x)>M_t(0)\) pro \(0<x<\tfrac12\)
      zde není dokázána.
\item Že společné jméno je komunikace z budoucnosti do minulosti,
      nebo že konečný záznam čerpá nekonečnou fyzikální energii.
\end{enumerate}

Věta~\ref{thm:factors} je uzavřená.
Říká, kdy se vlastnost smí jmenovat stejně jako odečet.
Nic víc.

\begin{thebibliography}{9}
\bibitem{poincare}
H.~Poincaré.
\newblock \emph{Science et méthode}.
\newblock Flammarion, Paris, 1908.
\newblock Věta o stejném jménu pro různé věci se tu bere jako čtení,
          ne jako doslovný zdroj definic.

\bibitem{lean}
J.~Mikulík.
\newblock \emph{finite-Shannon-theorem}, soubor \texttt{CausalSeed/GoldenSphere.lean}.
\newblock Commit \texttt{2e6b8e5}, 17.~listopadu 2026.
\newblock \url{https://github.com/jendamikulik/finite-Shannon-theorem}
\end{thebibliography}

\end{document}
